% =========================================================
% BÀI 6: BỀ MẶT HÀM CHI PHÍ VÀ QUỸ ĐẠO TỐI ƯU
% =========================================================
\section{Cost function landscape and optimization trajectory}

\subsection{Mục tiêu}
Phần này nhằm mục đích trực quan hóa hình dáng của hàm chi phí \(J(\theta_0, \theta_1)\) trong không gian 2 chiều và 3 chiều, đồng thời minh họa quỹ đạo cập nhật tham số (trajectory) của thuật toán Gradient Descent (GD) tự cài đặt từ điểm khởi tạo cho đến khi đạt điểm cực tiểu toàn cục.

\subsection{Kết quả thực hiện}

% Lưu ý: Bạn nhớ lưu bức ảnh 3D + 2D vào thư mục figures/ với tên bai_06_landscape.png nhé
\chenhinh{figures/06.jpg}
{Đồ thị 3D bề mặt hàm chi phí và Đồ thị 2D đường đồng mức (Contour)}
{fig:cost_landscape}

\subsection{Phân tích chi tiết}

\begin{itemize}
    \item \textbf{Không gian hàm chi phí:} Qua đồ thị 3D, ta có thể thấy hàm chi phí \(J(\boldsymbol{\theta})\) có dạng lòng chén (Paraboloid lồi). Đặc tính này đảm bảo rằng bài toán Hồi quy tuyến tính đơn biến với hàm Cost MSE luôn tồn tại một cực tiểu toàn cục duy nhất, do đó thuật toán Gradient Descent sẽ không bao giờ bị kẹt ở cực tiểu địa phương.
    
    \item \textbf{Quỹ đạo tối ưu hóa (GD Trajectory):} Đường màu đỏ thể hiện quá trình cập nhật tham số với \texttt{learning\_rate = 0.01} qua \(5000\) vòng lặp. Xuất phát từ điểm \(\boldsymbol{\theta} = [0,0]\) (nơi có giá trị Cost rất cao, khoảng \(92.38\)), quỹ đạo di chuyển mượt mà, men theo lòng chén và hội tụ về tâm một cách nhanh chóng.
    
    \item \textbf{Tính chính xác của mô hình:} Dấu X màu vàng là nghiệm giải tích (True Minimum) được tính từ phương trình chuẩn (Normal Equation) với giá trị Cost tối ưu là \(14.2574\). Điểm kết thúc của quỹ đạo Gradient Descent đã đè khít lên dấu X vàng. Điều này chứng minh thuật toán GD tự viết đã hoạt động chính xác tuyệt đối, tìm ra bộ tham số tương đương với nghiệm toán học.
    
    \item \textbf{Tính chất vuông góc:} Trên đồ thị 2D Contour, các đường đồng mức (Contour lines) thể hiện các mức Cost khác nhau. Quỹ đạo màu đỏ luôn cắt vuông góc (\(90^\circ\)) với các đường đồng mức. Điều này hoàn toàn khớp với lý thuyết: vector Gradient \(\nabla J(\boldsymbol{\theta})\) luôn vuông góc với đường đồng mức của hàm số tại điểm đó. 
    
    \textit{(Lưu ý: Đồ thị 2D bị kéo dẹt theo trục hoành do chênh lệch thang đo giữa \(\theta_0\) và \(\theta_1\), sử dụng \texttt{set\_aspect('equal')} để thấy rõ góc vuông).}
\end{itemize}